\section{RQ4--RQ5: Native Data and Agentic Episodes}
\label{sec:native}

The controlled catalogs isolate the effect of $N$, but they are artificial. We therefore repeat the measurements on data that needs no construction at all.

\begin{table}[t]
  \centering
  \caption{Native BFCL data (\%). Left: accuracy on \texttt{live\_multiple}, real user queries with the 2--37 tools that BFCL ships with them, by native tool count (25 queries per bucket); $^\dagger$lenient. Right: abstention rate on \texttt{live\_irrelevance}, where no tool fits and the correct behaviour is to make no call (25 per bucket).}
  \label{tab:native}
  \small
  \setlength{\tabcolsep}{4pt}
  \resizebox{\linewidth}{!}{\begin{tabular}{lccccccc|ccccc}
\toprule
 & \multicolumn{7}{c|}{\texttt{live\_multiple}: accuracy by native $N$} & \multicolumn{5}{c}{\texttt{live\_irrelevance}: abstention} \\
Model & 2 & 3 & 4 & 5 & 6+ & all & all$^\dagger$ & 1 & 2-3 & 4-5 & 6+ & all \\
\midrule
Qwen2.5-1.5B & 40.0 & 56.0 & 36.0 & 48.0 & 60.0 & 48.0 & 48.0 & 76.0 & 96.0 & 80.0 & 68.0 & 80.0 \\
Qwen2.5-3B & 68.0 & 84.0 & 68.0 & 68.0 & 60.0 & 69.6 & 69.6 & 64.0 & 80.0 & 56.0 & 84.0 & 71.0 \\
Llama-3.2-3B & 44.0 & 68.0 & 24.0 & 56.0 & 44.0 & 47.2 & 68.0 & 8.0 & 4.0 & 4.0 & 16.0 & 8.0 \\
Qwen3.5-2B & 56.0 & 72.0 & 68.0 & 76.0 & 76.0 & 69.6 & 69.6 & 56.0 & 56.0 & 52.0 & 80.0 & 61.0 \\
\bottomrule
\end{tabular}}
\end{table}

\paragraph{RQ4: real user queries.}
On \texttt{live\_multiple} (\cref{tab:native}), overall accuracy is \NatAccQwenOneFive\% (Qwen2.5-1.5B), \NatAccQwenThree\% (Qwen2.5-3B), \NatAccQwenThreeFive\% (Qwen3.5-2B) and \NatLenLlama\% for Llama-3.2-3B under lenient scoring (\NatAccLlama\% strict).
Accuracy shows no clear trend across native tool counts. This agrees with RQ1, because real BFCL queries mostly come with 2--10 tools, the range where the controlled curves are still nearly flat. With 25 queries per bucket, differences of up to about $\pm$19 points are noise.
As in RQ2, choosing the wrong tool is rare (\NatWrongToolLlama--\NatWrongToolQwenThree\% of queries). Argument values are the main difficulty.

\paragraph{Knowing when not to call.}
\texttt{live\_irrelevance} shows a failure mode that the controlled catalogs cannot.
The Qwen models correctly abstain in \NatAbstQwenThreeFive--\NatAbstQwenOneFive\% of cases, but \textbf{Llama-3.2-3B calls a tool \NatAbstInvLlama\% of the time when none applies} (abstention \NatAbstLlama\%). It inserts arguments that the user never gave.
More tools do not make abstention worse for any model in our sample. If anything, abstention is slightly higher in the largest bucket.
For an agent that executes tool calls, over-calling is arguably more dangerous than a missing call.

\begin{table}[t]
  \centering
  \caption{RQ5: 30 randomly sampled \texttt{multi\_turn\_base} episodes in BFCL's executable environments, scored by the official checker. \emph{Calls / GT calls}: executed calls relative to the ground truth. \emph{Dropped calls}: agent steps whose raw tokens contain a call that the server silently discarded. \emph{Forced quit}: share of episodes that hit the 20-step limit of a turn.}
  \label{tab:multiturn}
  \small
  \setlength{\tabcolsep}{4pt}
  \resizebox{\linewidth}{!}{\begin{tabular}{lcccccc}
\toprule
Model & Success (95\% CI) & Calls / GT calls & Exec.\ errors / ep. & Dropped calls & Forced quit & Top failure \\
\midrule
Qwen2.5-1.5B & 6.7 (1.8--21.3) & 0.50 & 0.9 & 5.9\% of steps & 0.0 & \small{empty turn model response} \\
Qwen2.5-3B & 23.3 (11.8--40.9) & 0.84 & 1.4 & 6.6\% of steps & 0.0 & \small{instance state mismatch} \\
Qwen3.5-2B & 20.0 (9.5--37.3) & 1.42 & 2.3 & 0.0\% of steps & 10.0 & \small{instance state mismatch} \\
\bottomrule
\end{tabular}}
\end{table}

\paragraph{RQ5: multi-turn agents.}
Stateful, multi-step episodes (\cref{tab:multiturn}; 9--42 tools, about 4 user turns each) are far harder than single calls.
Qwen2.5-3B completes \MtSuccQwenThree\% of episodes (95\% CI \MtLoQwenThree--\MtHiQwenThree), Qwen3.5-2B \MtSuccQwenThreeFive\%, and Qwen2.5-1.5B only \MtSuccQwenOneFive\%.
The models fail in different ways.
Qwen2.5-1.5B \emph{under-acts}: it makes only \MtCallsQwenOneFive{} calls per ground-truth call and often ends a turn without any call. \MtDroppedQwenOneFive\% of its steps contain a call that the server dropped because the JSON was invalid.
Qwen3.5-2B \emph{over-acts}: \MtCallsQwenThreeFive{} calls per ground-truth call, \MtExecErrQwenThreeFive{} execution errors per episode, and some turns that loop until the step limit.
Qwen2.5-3B lies in between; its typical failure is a final environment state that differs from the ground truth.
Two Qwen3.5-2B episodes first ended with an internal server error from Ollama. We re-ran both: one now completes (and fails for model reasons), while the other hits the same error again after 24 calls, probably because the ever-growing dialogue exhausts the context window (unverified). It is counted as a failure. Llama-3.2-3B was excluded from RQ5 because its 16k-token episodes exceeded the system memory of our laptop (\cref{sec:limitations}).
A typical failure is shown in \cref{app:example}. After an execution error (\texttt{lockDoors() missing 1 required positional argument: 'door'}), the 1.5B model gives up. In the next turn it misreads a tool result and reports tire pressures of 32--35\,psi as ``within the safe range of 40\,psi or higher''.

\paragraph{RQ3b: retrieval over a natively large corpus.}
On BFCL's 64-tool catalogs, the embedder almost always ranks the gold tool in the top 3 (\cref{sec:rq3}). ToolRet's native corpus is much harder: 44{,}453 tools aggregated from 35 tool-use datasets (\cref{tab:toolret}).
Retrieving against the full corpus, BGE-M3 reaches nDCG@10 of \TrNdcgBge{} and Qwen3-Embedding-0.6B \TrNdcgQwenGen{}. With ToolRet's per-query task instruction, Qwen3-Embedding reaches \TrNdcgQwenTask, but still finds only \TrRecallQwenTaskKTen\% of the relevant tools in its top 10.
Small LLMs therefore need a short tool list, and with realistic tool libraries a short list will often miss the right tool. In such deployments, the \emph{retriever}, not the LLM, becomes the bottleneck. This agrees with \citet{shi2025toolret}.

\begin{table}[t]
  \centering
  \caption{RQ3b: retrieval on ToolRet's full native corpus (44{,}453 tools; 600 queries, 200 per category). Recall is the share of a query's relevant tools that are retrieved, averaged over queries.}
  \label{tab:toolret}
  \small
  \begin{tabular}{lcccc|cc}
\toprule
Retriever & \multicolumn{4}{c|}{nDCG@10} & \multicolumn{2}{c}{Recall (all)} \\
 & web & code & custom. & all & @10 & @50 \\
\midrule
bge-m3 & 33.7 & 25.4 & 22.3 & 27.1 & 35.4 & 48.1 \\
qwen3-emb (generic instr.) & 44.6 & 43.3 & 32.2 & 40.0 & 50.9 & 66.9 \\
qwen3-emb (task instr.) & 48.1 & 58.0 & 40.0 & 48.7 & 58.7 & 73.8 \\
\bottomrule
\end{tabular}
\end{table}
